\subsection{Interactions}
\label{sec:interactions}

Random forests are particularly adept at capturing interactions between
features and are
for that reason often utilized when the features are believed to interact.
Here, we examine how variable importances are allocated when a feature is
itself an interaction between other features.

Consider a pure two-way parity interaction with no marginal effects:
\[
  y = c/10 + \big((a + b) \bmod 2\big), \qquad a,b,c \in \{0,1\}, \quad n =
  10000.
\]
Neither $a$ nor $b$ carries information about $y$ on its own; only their joint
parity does. $c$ is the only variable with an independent effect. We fit a
random forest and a linear regression model against this data.

Main-effects \textbf{OLS} sees only $c$: the fitted coefficients on $a$ and $b$
are $0.000$ and $0.000$ against $0.101$ for $c$, at $R^2 = 0.01$.
For binary inputs $(a+b) \bmod 2 = a + b - 2ab$ is bilinear, so adding an
$a{\times}b$ term fits the target exactly ($R^2 = 1.000$).

When the interaction is captured in the Random Forest, the two interacting
variables $a$ and $b$
are allocated most of the variable importance. The size of their share follows
from the variance decomposition. Impurity importance sums the reduction in node
sum of squares over splits, and
trees grown to purity remove $n\,\mathrm{Var}(y)$ between them, so the
importance
divides as the variance does. The target is additive in two independent pieces,
$\mathrm{Var}(c/10) = 0.0025$ and $\mathrm{Var}\big((a+b) \bmod 2\big) = 0.25$:
the
coefficient $1/10$ gives $c$ one hundredth of the pair's variance, so $c$ gets
importance
$1\%$ and the pair $99\%$. Note that this is exactly the explained
variance ($R^2 = 0.01$) of main effects OLS.
At $\code{mtry} = 3$, the measured total is $2525.2 = n\,\mathrm{Var}(y)$ and
$c$ receives $1.02\%$
of it. Over seeds $1$--$10$ at the default $\code{mtry} = 1$, $a$ and $b$
average
$451$ (SD $29$) against $15$ for $c$, and its share rises to $1.6\%$, because
$\code{mtry} = 1$ truncates trees and the loss falls harder on the pair, whose
credit needs both variables on one path: a node whose single candidate is
already constant becomes terminal, so paths are cut before both $a$ and $b$
have been split on. Figure~\ref{fig:interaction} shows the contrast
against main-effects OLS.

Since $a$ and $b$ are exchangeable by construction, any difference between their
importances is sampling noise, and within a single seed they spread by
$2$--$20\%$ (median $10\%$), against $1$--$6\%$ when the parity is replaced by
equal \emph{marginal} effects for $a$ and $b$ ($y = c/10 + (a + b)/2$).
So only the $a,b \gg c$ gap is informative here.

The interaction spread is the wider one because credit for an interacting
variable is
conditional on tree topology. Splitting on $a$ or $b$ alone reduces no variance;
once a tree has split on one of the pair, the
other separates the parity completely. Each tree therefore hands the pair's
entire contribution to whichever variable it splits last.
At $\code{mtry} = 1$ that division between $a$ and $b$ is made by the
candidate draw rather than by the data,
so the number of trees crediting each variable is binomial and the imbalance
averages
away as trees accumulate.
At
$\code{mtry} = p$ both variables compete at each node and the one with the
larger
empirical decrease wins:
with bagging disabled the pair's credit goes entirely to one variable, and
bootstrap resampling only partly decorrelates that choice.
Conditional on the
training sample the surviving imbalance therefore becomes a bias rather than
noise:
every tree re-derives the same allocation from the same data, so averaging
converges to that allocation instead of to an even split. Refitting each of six
datasets with
eight forest seeds at $\code{mtry} = p$ confirms this: all eight land on the
same side of an even split for every dataset, and from $500$ to $4000$ trees the
forest-to-forest SD falls from $0.02$ to $0.005$ while the mean does not move.
The size of the imbalance is the sample's: the $a$-vs-$b$ spread runs from
$3\%$ to $119\%$ across the six.

Breiman introduced the per-node
random draw of $F$ candidate inputs (what the parameter \code{mtry} implements)
because
``the randomness injected has to minimize the correlation $\bar\rho$ while
maintaining strength'' \citep[\S3; Thm.~11.2 for regression]{breiman2001},
$\bar\rho$ being the correlation between the predictions of pairs of trees.
At $\code{mtry} = p$ none of that randomness is injected. The $\code{mtry}$
setting that results in the best OOB generalization (OOB $R^2$ $1.00$ at
$\code{mtry} = p$ against $0.59$ at the default) is therefore also the one whose
division of
importance between $a$ and $b$ is least trustworthy.

AIR was not designed to remove \emph{this bias}, but it has an incidental
effect through tree decorrelation.
At $\code{mtry} = 1$ the two importance measures are indistinguishable, whereas
AIR shows improvement
at $\code{mtry} = p$ (median $a$-vs-$b$ spread $13.3\%$ against $29.1\%$).
This comes
from its permuted pseudo-variables enlarging the candidate pool to $2p$, which
lowers $\bar\rho$ in the same vein as a smaller \code{mtry} would, degrading
also
OOB $R^2$ to $0.946$.

The same allocation question can be put to the two measures most often used to
corroborate an impurity ranking, and which this paper takes as its foils:
\textbf{permutation importance} (OOB Breiman-style, as returned by
\code{ranger}) and \textbf{mean absolute TreeSHAP} ($\overline{|\varphi_j|}$,
the path-dependent \emph{observational} variant). Both confirm the $a,b \gg c$
reading, so switching measure does not rescue the interpretation, and both give
$a$ and $b$ equal scores. They part company on how the pair's shared credit is
divided. Permutation importance charges each member the \emph{entire}
interaction rather than a share of it --- scrambling either variable destroys
the whole parity --- so the pair's importances sum to twice the interaction
variance, while mean $|\varphi|$ partitions the credit exactly. Each is
denominated in units of its own, so neither is an effect size.
Appendix~\ref{app:interactions} derives both; Table~\ref{tab:interactions}
collects the allocations.

\begin{table}[htb]
\centering
\caption{The four allocations on the seed-42 parity dataset, ratios normalised
to $c$ (the only variable with a marginal effect).}
\label{tab:interactions}
\begin{tabular}{llccc}
\toprule
\code{mtry} & measure & $a : b : c$ (to $c$) & $c$ share & $a/(a{+}b)$ \\
\midrule
1 (default) & MDI              & $28 : 23 : 1$    & $1.9\%$  & $0.545$ \\
            & AIR              & $21 : 22 : 1$    & $2.3\%$  & $0.496$ \\
            & permutation      & $53 : 53 : 1$    & $0.9\%$  & $0.501$ \\
            & mean $|\varphi|$ & $2.7 : 2.7 : 1$  & $15.8\%$ & $0.500$ \\
\midrule
3 ($= p$)   & MDI              & $44 : 52 : 1$    & $1.0\%$  & $0.459$ \\
            & AIR              & $40 : 39 : 1$    & $1.2\%$  & $0.503$ \\
            & permutation      & $100 : 100 : 1$  & $0.5\%$  & $0.500$ \\
            & mean $|\varphi|$ & $5 : 5 : 1$      & $9.1\%$  & $0.500$ \\
\bottomrule
\end{tabular}
\end{table}

\paragraph{Practical implications.} A variable at the top of an impurity
ranking may have no marginal effect at all and owe its place to an interaction,
and switching measure does not rescue the reading: permutation importance and
mean $|\varphi|$ rank the interacting pair as far above $c$ as MDI does, while
within an interacting pair, the order is not to be read --- the division of
their shared credit is candidate-draw noise at the default \code{mtry} and a
sample-fixed bias at $\code{mtry} = p$ under MDI, and permutation importance
additionally charges each member the full joint effect rather than a share of
it.
What distinguishes the two cases is the one-way
partial-dependence plot (PDP) under a tall bar: here it is flat, and the parity
appears only in the joint PDP over $(a,b)$.

\begin{figure}[htb]
\centering
\includegraphics[width=0.85\textwidth]{fig2_interaction.png}
\caption{Same data, two views. Main-effects OLS assigns $a$ and $b$ essentially
zero coefficient, seeing only $c$; RF impurity importance ranks the two
interacting variables far above it. The ordering between $a$ and $b$ is sampling
noise --- they are exchangeable by construction.}
\label{fig:interaction}
\end{figure}
